\section*{Hoofdstuk \ref{ch:g9:kinetic-energy-safety} --- Kinetische energie en verkeersveiligheid}

\begin{solution}{exo:g9:kinetic-energy-safety:1}
De \omterm{def:g9:kinetic-energy-safety:joule}{joule} (\unit{J}): een paar \omterm{def:g9:kinetic-energy-safety:joule}{joule} tillen dit boek
naar zijn plank; het hart geeft er ongeveer één per slag uit. $E_k =
\frac12 m v^2$ --- $m$ in \omterm{def:g3:measuring-mass:units}{kilogram}, $v$ in
\omterm{def:g2:measuring-length:units}{meter} per \omterm{def:g2:measuring-time:second}{seconde}, $E_k$ in \omterm{def:g9:kinetic-energy-safety:joule}{joule}.
\end{solution}

\begin{solution}{exo:g9:kinetic-energy-safety:2}
Haas: $0.5 \times 2 \times 100 = \qty{100}{J}$. Speler: $0.5 \times 80
\times 25 = \qty{1000}{J}$.
\end{solution}

\begin{solution}{exo:g9:kinetic-energy-safety:3}
De helft meer \omterm{def:g3:measuring-mass:mass}{massa}: factor $1.5$. De helft meer
\omterm{def:g7:motion-average-speed:formula}{snelheid}: factor $1.5^2 = 2.25$. De snelheidsknop wint --- dat
doet hij altijd, want hij draagt het kwadraat.
\end{solution}

\begin{solution}{exo:g9:kinetic-energy-safety:4}
\omterm{prop:g9:kinetic-energy-safety:stopping}{Reactieafstand} --- de weg van één \omterm{def:g2:measuring-time:second}{seconde}, evenredig groeiend
met $v$; \omterm{prop:g9:kinetic-energy-safety:stopping}{remweg} --- de pensionering van $E_k$, groeiend met
$v^2$.
\end{solution}

\begin{solution}{exo:g9:kinetic-energy-safety:5}
Bij \qty{50}{km/h}: $13$ van de \qty{27}{m} --- ongeveer de helft. Bij
\qty{130}{km/h}: $85$ van de \qty{121}{m} --- zo'n zeventig procent.
Het gekwadrateerde onderdeel groeit het evenredige voorbij: bij
\omterm{def:g7:motion-average-speed:formula}{snelheid} verslindt het remmen het totaal.
\end{solution}

\begin{solution}{exo:g9:kinetic-energy-safety:6}
Naar \omterm{def:g5:heat-and-insulation:heat}{warmte} in de remschijven --- de amberkleurige gloed na
afdalingen in de bergen, de geur van hete remmen.
\omterm{def:g5:everyday-energy:energy}{Energie} wordt doorgegeven, nooit uitgewist: de ketenpropositie
uit het energiehoofdstuk van je kindertijd, nog altijd aan de
macht.
\end{solution}

\begin{solution}{exo:g9:kinetic-energy-safety:7}
$36 \div 3.6 = \qty{10}{m}$; dan \qty{20}{m}; dan \qty{30}{m} ---
marcherend evenredig met de \omterm{def:g7:motion-average-speed:formula}{snelheid}, zoals een stuk van één
\omterm{def:g2:measuring-time:second}{seconde} wel moet.
\end{solution}

\begin{solution}{exo:g9:kinetic-energy-safety:8}
Je $E_k$ ligt vast door je \omterm{def:g3:measuring-mass:mass}{massa} en de
\omterm{def:g7:motion-average-speed:formula}{snelheid} van de auto --- geen enkele band verandert daar iets
aan. De gordel beschaaft de \emph{pensionering}: hij stopt het lichaam
over een langere afstand en tijd, zachtjes, in plaats van tegen het
dashboard, in één klap.
\end{solution}

\begin{solution}{exo:g9:kinetic-energy-safety:9}
Het remmen schaalt met het kwadraat van de \omterm{def:g7:motion-average-speed:formula}{snelheid}: bij
\qty{60}{km/h} (dubbel), $8 \times 4 = \qty{32}{m}$; bij \qty{90}{km/h}
(drievoudig), $8 \times 9 = \qty{72}{m}$.
\end{solution}

\begin{solution}{exo:g9:kinetic-energy-safety:10}
Bij \qty{30}{km/h}: reactie $\approx \qty{8.3}{m}$, remmen $13 \times
(30/50)^2 \approx \qty{4.7}{m}$ --- stoppen in ongeveer \qty{13}{m},
tegen \qty{27}{m} bij vijftig: minder dan de halve weg. Energieën:
$(30/50)^2 = 0.36$ --- nauwelijks een derde van de botsenergie. Allebei
overtuigen ze; de posterklare regel is meestal die van de weg --- ``bij
30 sta je stil vóór het kind; bij 50 bereik je het nog in beweging''.
\end{solution}

\begin{solution}{exo:g9:kinetic-energy-safety:11}
Verdubbel alleen het remaandeel --- regen rekt de pensionering, niet het
brein: bij \qty{50}{km/h} $14 + 26 = \qty{40}{m}$; bij \qty{90}{km/h}
$25 + 80 = \qty{105}{m}$. De reactieseconde is droog en nat hetzelfde.
\end{solution}

\begin{solution}{exo:g9:kinetic-energy-safety:12}
Vrachtwagen: $v = 25$ \unit{m/s}; $E_k = 0.5 \times 20000 \times 625 =
\qty{6.25e6}{J}$. Bijpassende auto: $0.5 \times 1000 \times v^2 =
\num{6.25e6}$ geeft $v^2 = 12500$, $v \approx \qty{112}{m/s} \approx
\qty{400}{km/h}$ --- geen enkele personenauto komt in de buurt. Een
vrachtwagen op snelwegsnelheid draagt raceauto-energieën met
vrachtwagenremmen: als die op een lange afdaling oververhitten en
falen, kan alleen een grindbaan megajoules veilig met pensioen sturen.
Auto's hebben er nooit een nodig; hun energieën sterven in gewone
remmen.
\end{solution}

\begin{solution}{pb:g9:kinetic-energy-safety:1}
\textbf{1.} $50 \div 3.6 = \qty{13.9}{m/s}$; $60 \div 3.6 =
\qty{16.7}{m/s}$.
\textbf{2.} \qty{13.9}{m} en \qty{16.7}{m}.
\textbf{3.} \qty{13}{m}; bij zestig $13 \times (60/50)^2 = 13 \times
1.44 \approx \qty{18.7}{m}$.
\textbf{4.} Stoppen: $26.9$ tegen \qty{35.4}{m} --- het gat in de
poster: ongeveer \emph{achtenhalve \omterm{def:g2:measuring-length:units}{meter}} meer, twee
autolengtes voorbij de plek waar de vijftigrijder al stilstaat.
\textbf{5.} Bij vijftig: ongeveer \qty{97}{kJ}. Bij negentig ($v =
\qty{25}{m/s}$): $0.5 \times 1000 \times 625 = \qty{312}{kJ}$.
\textbf{6.} $97 \div 9.8 \approx \qty{10}{m}$; $312 \div 9.8 \approx
\qty{32}{m}$.
\textbf{7.} Tien \omterm{def:g2:measuring-length:units}{meter}: een huis van drie verdiepingen;
tweeëndertig: een flat van tien. Ontwerp: ``Een botsing bij 50 is een
val van de derde verdieping. Bij 90 van de tiende. Kies je raam.''
\textbf{8.} De hoogtes \omterm{def:g6:measurement-in-science:measuring}{meten} eerlijk de \omterm{def:g5:everyday-energy:energy}{energie} die
met pensioen moet --- aan dat getal verandert geen enkele techniek iets.
Kreukelzones en gordels veranderen de \emph{manier} van pensionering:
\omterm{def:g2:measuring-length:units}{meters} zacht stoppen in plaats van \omterm{def:g2:measuring-length:units}{centimeters} wreed
stoppen --- en daarom overdrijft de valvergelijking het letsel in een
moderne auto, en is de \omterm{def:g5:everyday-energy:energy}{energie} die ze afbeeldt toch echt.
\textbf{9.} \qty{25}{m} blind per \omterm{def:g2:measuring-time:second}{seconde}; een blik van twee \omterm{def:g2:measuring-time:second}{seconden}
\qty{50}{m} --- een half veld voordat de reactieseconde überhaupt
begint.
\textbf{10.} $65 + 25 = \qty{90}{m}$ --- de vermoeide bestuurder legt
anderhalve buslengte bij het totaal van de oplettende bestuurder.
\textbf{11.} Het meest mis bij \emph{lage} snelheden: daar is het
gekwadrateerde remaandeel klein en vormt het reactiestuk het grootste
deel van de stopafstand --- in de stad \emph{is} die \omterm{def:g2:measuring-time:second}{seconde} het gevaar.
(Bij hoge \omterm{def:g7:motion-average-speed:formula}{snelheid} overheerst het remmen --- maar \qty{25}{m}
per \omterm{def:g2:measuring-time:second}{seconde} is ook daar geen kleinigheid.)
\textbf{12.} Bijvoorbeeld: ``Tien \omterm{def:g7:motion-average-speed:formula}{snelheid} erbij, tien
\omterm{def:g2:measuring-length:units}{meter} stoppen erbij --- de \omterm{def:g7:motion-average-speed:formula}{snelheid} telt twee keer
mee, één keer gekwadrateerd. Een botsing is een val: de
\omterm{def:g5:everyday-energy:energy}{energie} groeit met het kwadraat van de
\omterm{def:g7:motion-average-speed:formula}{snelheid}. Eén \omterm{def:g2:measuring-time:second}{seconde} blik is vijfentwintig
\omterm{def:g2:measuring-length:units}{meter} blindheid --- aandacht is de kortste
\omterm{prop:g9:kinetic-energy-safety:stopping}{remweg} van allemaal.''
\end{solution}
